\chapter{Complexity — Meyer Sound}
%%% generated by the blueprint pipeline from TCSlib.Complexity.CircuitComplexity.MeyerSound
\section{Overview}

Declarations of the Lean module \texttt{TCSlib.Complexity.CircuitComplexity.MeyerSound}.

\section{Declarations}

\begin{definition}[Time-variable encoding]
\lean{Complexity.Meyer.tvOf}
\label{Complexity.Meyer.tvOf}
\leanok
\uses{Complexity.Meyer.TW}
Let $W,t\in\mathbb{N}$. Define a map from four symbols to natural numbers by sending the
current-time symbol to $t$, its successor symbol to $t+1$, the zero symbol to $0$, and
the final symbol to $2^W-1$.
\end{definition}

\begin{definition}[Address-variable encoding]
\lean{Complexity.Meyer.avOf}
\label{Complexity.Meyer.avOf}
\leanok
\uses{Complexity.Meyer.AW}
Let $a,j\in\mathbb{N}$. Define a map from six symbols to natural numbers by sending the
current-address symbol to $a$, its successor symbol to $a+1$, the zero symbol to $0$,
the one symbol to $1$, the address-$j$ symbol to $j$, and its successor symbol to $j+1$.
\end{definition}

\begin{definition}[Boolean valuation of atoms]
\lean{Complexity.Meyer.NumVal}
\label{Complexity.Meyer.NumVal}
\leanok
\uses{Complexity.Meyer.Atom, Complexity.Meyer.Sel, Complexity.Meyer.avOf, Complexity.Meyer.tvOf}
Let $M$ be a finite Turing machine with Boolean tape alphabet, let $\mathrm{Ans}$ assign
a Boolean value to each selector and pair of natural numbers, and let
$W,n,t,a,j\in\mathbb{N}$ and $xb\in\{\mathrm{false},\mathrm{true}\}$. Define the Boolean
value of each atom as follows: an answer atom is evaluated by applying $\mathrm{Ans}$ to
its selector with the encoded time and address; the length atom is true; the time and
address not-at-their-maximum atoms assert respectively $t+1<2^W$ and $a+1<2^W$; the zero
and index atoms assert respectively $a=0$, $j=0$, $j<n$, $j\le n$, and $j=n+1$; and the
input-bit atom has value $xb$.
\end{definition}

\begin{definition}[Symbol-value encoding]
\lean{Complexity.Meyer.sv}
\label{Complexity.Meyer.sv}
\leanok
Let $M$ be a finite Turing machine with Boolean tape alphabet, and let $\mathrm{Ans}$
assign a Boolean value to each selector and pair of natural numbers. For an optional
tape position $\tau$, a Boolean sign $s$, and natural numbers $t,a$, define the encoded
symbol value to be absent if the answer for sign $\mathrm{true}$ is true, and otherwise
to be the answer for sign $\mathrm{false}$ at $(\tau,t,a)$.
\end{definition}

\begin{definition}[Signed-coordinate symbol encoding]
\lean{Complexity.Meyer.cX}
\label{Complexity.Meyer.cX}
\leanok
\uses{Complexity.Meyer.sv}
Let $M$ be a finite Turing machine with Boolean tape alphabet, and let $\mathrm{Ans}$
assign a Boolean value to each selector and pair of natural numbers. For an optional
tape position $\tau$, a natural number $t$, and an integer $d$, define the encoded
symbol value using time $t$, address $|d|$, and sign equal to true exactly when $d<0$.
\end{definition}

\begin{theorem}[Symbolic valuation agrees with table answers]
\lean{Complexity.Meyer.symvA_NumVal}
\label{Complexity.Meyer.symvA_NumVal}
\leanok
\uses{Complexity.Meyer.AW, Complexity.Meyer.NumVal, Complexity.Meyer.TW, Complexity.Meyer.avOf, Complexity.Meyer.sv, Complexity.Meyer.symvA, Complexity.Meyer.tvOf}
\difficulty{1}
Let $M$ be a finite-tape machine, let $\mathrm{Ans}$ be a Boolean table of selector
answers, and choose natural numbers $W,n,t,a,j$ and a Boolean $b$. For every tape
selector $\tau$, Boolean sign $s$, time offset label $T$, and coordinate label $A$,
evaluating the symbolic-answer atom under the numerical valuation determined by these
parameters gives the same optional Boolean value as the table query selected by $\tau,s$
at the time coordinate represented by $T$ and the position coordinate represented by
$A$.
\end{theorem}

\begin{theorem}[State Answers Under Numerical Encoding]
\lean{Complexity.Meyer.stA_NumVal}
\label{Complexity.Meyer.stA_NumVal}
\leanok
\uses{Complexity.Meyer.NumVal, Complexity.Meyer.TW, Complexity.Meyer.stA, Complexity.Meyer.tvOf, Complexity.TimeHierarchy.OutReg}
\difficulty{1}
Let $M$ be a finite Turing machine with Boolean alphabet, let $\mathrm{Ans}$ be a
Boolean-valued function on selectors and two natural-number indices, and choose natural
numbers $W,n,t,a,j$ and a Boolean $xb$. For every state-and-output-register pair
$\sigma$ and time coordinate $T$, the Boolean value assigned to the state-answer atom
for $\sigma$ at coordinates $(T,0)$ by the numerical encoding is
$\mathrm{Ans}(\mathrm{st}(\sigma),\mathrm{time}(T),0)$, where $\mathrm{time}(T)$ is $t$,
$t+1$, $0$, or $2^W-1$ according as $T$ denotes the current, next, initial, or last
time.
\end{theorem}

\begin{theorem}[Value of an Answer Atom]
\lean{Complexity.Meyer.NumVal_ans}
\label{Complexity.Meyer.NumVal_ans}
\leanok
\uses{Complexity.Meyer.AW, Complexity.Meyer.NumVal, Complexity.Meyer.Sel, Complexity.Meyer.TW, Complexity.Meyer.avOf, Complexity.Meyer.tvOf}
\difficulty{1}
Let $M$ be a finite-tape Turing machine with Boolean tape symbols, let $\mathrm{Ans}$
assign a Boolean to each table selector and pair of natural-number indices, and let
$W,n,t,a,j$ be natural numbers and $xb$ a Boolean. For any selector $s$, time tag $T$,
and address tag $A$, the numerical valuation of the atom asserting the answer at
$(s,T,A)$ is $\mathrm{Ans}(s,t_T,a_A)$, where $t_T$ is respectively $t$, $t+1$, $0$, or
$2^W-1$ for the four time tags, and $a_A$ is respectively $a$, $a+1$, $0$, $1$, $j$, or
$j+1$ for the six address tags.
\end{theorem}

\begin{theorem}[Decoding a Boolean symbol]
\lean{Complexity.Meyer.decode_symBit}
\label{Complexity.Meyer.decode_symBit}
\leanok
\uses{Complexity.Meyer.symBit}
\difficulty{4}
For every optional Boolean symbol $v$, let the blank test return true exactly when $v$
is absent, and let the one test return true exactly when $v$ is the symbol true. The
optional symbol reconstructed as absent when the blank test is true, and otherwise as
the result of the one test, equals $v$.
\end{theorem}

\begin{definition}[Integer sign indicator]
\lean{Complexity.Meyer.sE}
\label{Complexity.Meyer.sE}
\leanok
For an integer $e$, define its sign indicator to be true exactly when $e<0$.
\end{definition}

\begin{definition}[Natural-number encoding of integers]
\lean{Complexity.Meyer.aE}
\label{Complexity.Meyer.aE}
\leanok
For an integer $e$, define its natural-number encoding to be $-e-1$ when $e<0$, and $e$
otherwise, with the resulting integer viewed as a natural number.
\end{definition}

\begin{theorem}[Signed integer reconstruction]
\lean{Complexity.Meyer.edge_c1}
\label{Complexity.Meyer.edge_c1}
\leanok
\uses{Complexity.Meyer.aE, Complexity.Meyer.sE, Complexity.Meyer.sgnZ}
\difficulty{4}
For every integer $e$, let its sign bit indicate whether it is negative and let its
unsigned index be $-e-1$ when it is negative and $e$ otherwise. Recovering an integer by
negating the index plus one in the negative case, and using the index otherwise, gives
$e$.
\end{theorem}

\begin{theorem}[Successor reconstruction from signed index]
\lean{Complexity.Meyer.edge_c2}
\label{Complexity.Meyer.edge_c2}
\leanok
\uses{Complexity.Meyer.aE, Complexity.Meyer.sE, Complexity.Meyer.sgnZ}
\difficulty{4}
For every integer $e$, let its sign bit indicate whether it is negative and let its
unsigned index be $-e-1$ when it is negative and $e$ otherwise. Recovering an integer by
negating the index when it is negative, and using the index plus one otherwise, gives
$e+1$.
\end{theorem}

\begin{theorem}[Negative endpoint indicator]
\lean{Complexity.Meyer.edge_E1}
\label{Complexity.Meyer.edge_E1}
\leanok
\uses{Complexity.Meyer.aE, Complexity.Meyer.sE}
\difficulty{4}
For every integer $e$, define its sign bit to indicate whether $e<0$, and define its
unsigned index to be $-e-1$ when $e<0$ and $e$ otherwise. The conjunction that $e$ is
negative and its unsigned index is zero is true exactly when $e+1=0$.
\end{theorem}

\begin{theorem}[Zero endpoint indicator]
\lean{Complexity.Meyer.edge_E0}
\label{Complexity.Meyer.edge_E0}
\leanok
\uses{Complexity.Meyer.aE, Complexity.Meyer.sE}
\difficulty{4}
For every integer $e$, define its sign bit to indicate whether $e<0$, and define its
unsigned index to be $-e-1$ when $e<0$ and $e$ otherwise. The conjunction that $e$ is
nonnegative and its unsigned index is zero is true exactly when $e=0$.
\end{theorem}

\begin{theorem}[Successor identity for signed integer encoding]
\lean{Complexity.Meyer.aE_add_one}
\label{Complexity.Meyer.aE_add_one}
\leanok
\uses{Complexity.Meyer.aE}
\difficulty{3}
Let $e$ be an integer, and define $a(e)$ by $a(e)=-e-1$ when $e<0$ and $a(e)=e$ when
$e\geq 0$. Then $a(e)+1$ equals $|e|$ when $e<0$, and equals $|e|+1$ when $e\geq 0$.
\end{theorem}

\begin{definition}[Validity of all clauses]
\lean{Complexity.Meyer.AllGood}
\label{Complexity.Meyer.AllGood}
\leanok
\uses{Complexity.Meyer.NumVal, Complexity.Meyer.goodF}
Let $M$ be a finite Turing machine with Boolean tape alphabet, let $\mathrm{Ans}$ assign
a Boolean value to each selector and pair of natural numbers, let $W\in\mathbb{N}$, let
$x$ be a finite Boolean list, and let
$xb:\mathbb{N}\to\{\mathrm{false},\mathrm{true}\}$. For every $t,a,j\in\mathbb{N}$ with
$t<2^W$, $a<2^W$, and $j\le |x|+1$, and for every clause $\pi$, the clause evaluates to
true under the good-clause evaluation using $\mathrm{Ans}$, the parameters
$W,|x|,t,a,j$, and input bit $xb(j)$.
\end{definition}

\begin{theorem}[Symmetric answers at zero]
\lean{Complexity.Meyer.sv_negZero}
\label{Complexity.Meyer.sv_negZero}
\leanok
\uses{Complexity.Meyer.AllGood, Complexity.Meyer.NumVal_ans, Complexity.Meyer.avOf, Complexity.Meyer.goodF, Complexity.Meyer.sv, Complexity.Meyer.tvOf}
\difficulty{4}
Let $M$ be a finite-tape Turing machine with Boolean symbols, let $x$ be a Boolean
input, and let $W$ be a natural number. Suppose a table of Boolean answers, indexed by
machine-state or tape-symbol queries and two natural-number coordinates, satisfies every
checker condition for all $t,a,j$ with $t,a<2^W$ and $j\le |x|+1$, using the input bit
supplied at position $j$. Then for every tape selector $\tau$, every $t<2^W$, and the
second coordinate $0$, the two symbol queries distinguished by the Boolean sign
parameter have equal encoded answers:
$\mathrm{sv}(\tau,\mathrm{true},t,0)=\mathrm{sv}(\tau,\mathrm{false},t,0)$.
\end{theorem}

\begin{theorem}[Signed coordinates for symbolic answers]
\lean{Complexity.Meyer.sv_sgn}
\label{Complexity.Meyer.sv_sgn}
\leanok
\uses{Complexity.Meyer.AllGood, Complexity.Meyer.cX, Complexity.Meyer.sgnZ, Complexity.Meyer.sv, Complexity.Meyer.sv_negZero}
\difficulty{4}
Let $M$ be a finite-tape Turing machine with Boolean symbols, let $x$ be a Boolean
input, and let $W$ be a natural number. Suppose a table of Boolean answers satisfies
every checker condition for all $t,a<2^W$ and $j\le |x|+1$, using the supplied input-bit
function. For every selector $\tau$, time $t<2^W$, Boolean sign $s$, and natural number
$a$, the symbolic answer at coordinates $(t,a)$ equals the answer at time $t$ and signed
integer coordinate $(-1)^s a$, where the sign convention is $(-1)^{\mathrm{false}}a=a$
and $(-1)^{\mathrm{true}}a=-a$.
\end{theorem}

\begin{theorem}[Zero-offset symbol query]
\lean{Complexity.Meyer.cX_zero}
\label{Complexity.Meyer.cX_zero}
\leanok
\uses{Complexity.Meyer.cX, Complexity.Meyer.sv}
\difficulty{1}
Let $M$ be a finite-tape Turing machine with Boolean tape symbols, and let $A$ assign a
Boolean answer to each state-or-symbol query and each pair of natural-number indices.
For any optional tape index $\tau$ and time $t$, the signed query at offset zero equals
the symbol query with the negative-sign flag set to false.
\end{theorem}

\begin{definition}[Correctness of the answer table]
\lean{Complexity.Meyer.Correct}
\label{Complexity.Meyer.Correct}
\leanok
\uses{Complexity.Meyer.RI, Complexity.Meyer.RW, Complexity.Meyer.SS, Complexity.Meyer.cX, Complexity.TimeHierarchy.OutReg, Turing.MultiTapeTM.initCfg, Turing.MultiTapeTM.runFrom}
Let $M$ be a finite Turing machine with Boolean tape symbols, let $x$ be a finite
Boolean input, let $W,t\in\mathbb{N}$, and let $\mathrm{Ans}$ be an answer function for
selectors of $M$. The predicate holds when, after $t$ steps from the initial
configuration on $x$, the answer function on every state and output-register pair agrees
with whether that pair is the machine’s actual state and output; for each work tape $h$
and offset $d\in\mathbb{Z}$ with $|d|+t\le 2^W-1$, the table’s work-tape value agrees
with the machine’s symbol at that offset from the tape head; and for every
$d\in\mathbb{Z}$ with $|d|\le |x|+1$, its input-tape value agrees with the machine’s
input symbol at that offset from the input head.
\end{definition}

\begin{theorem}[Correctness implies true claims]
\lean{Complexity.Meyer.claims_of_correct}
\label{Complexity.Meyer.claims_of_correct}
\leanok
\uses{Complexity.Meyer.Correct, Complexity.Meyer.NumVal, Complexity.Meyer.SS, Complexity.Meyer.avOf, Complexity.Meyer.cX_zero, Complexity.Meyer.claims, Complexity.Meyer.reads, Complexity.Meyer.stA_NumVal, Complexity.Meyer.symvA_NumVal, Complexity.Meyer.tvOf, Turing.MultiTapeTM.initCfg, Turing.MultiTapeTM.runFrom}
\difficulty{4}
Let $M$ be a finite-tape Turing machine with Boolean symbols, let $x$ be a finite
Boolean input, and let $A$ be a Boolean answer table for state and tape-symbol queries
indexed by two natural numbers. Fix a width $W$, time $t$, indices $a,j$, and a Boolean
$b$. Suppose that at time $t$, the table answers every state-and-output query exactly
according to the machine's configuration, and answers every work-tape query at signed
offset $d$ correctly whenever $|d|+t\le 2^W-1$, as well as every input query whenever
$|d|\le |x|+1$; also suppose $t<2^W$. Then the numerical valuation determined by
$A,W,|x|,t,a,j,b$ makes the state-and-read claims at time $t$ true for the machine's
actual configuration.
\end{theorem}

\begin{theorem}[Signed index agrees with shifted word]
\lean{Complexity.Meyer.xin_eq}
\label{Complexity.Meyer.xin_eq}
\leanok
\uses{Complexity.Meyer.xhat}
\difficulty{4}
Let $x$ be a Boolean word, let $xb$ assign a bit to each natural number, and suppose
$xb(j)=x_j$ for every $j<|x|$. For every integer $d$, the following expression equals
the symbol of $x$ shifted to position $1+d$: it is blank if $d<0$ and $d\ne0$, equals
$xb(|d|)$ if $|d|<|x|$ otherwise, and is blank if $|d|\ge |x|$ otherwise.
\end{theorem}

\begin{theorem}[Correctness at the Initial Time]
\lean{Complexity.Meyer.correct_zero}
\label{Complexity.Meyer.correct_zero}
\leanok
\uses{Complexity.Meyer.AllGood, Complexity.Meyer.Correct, Complexity.Meyer.NumVal, Complexity.Meyer.NumVal_ans, Complexity.Meyer.RI_init, Complexity.Meyer.RW_init, Complexity.Meyer.SS_init, Complexity.Meyer.avOf, Complexity.Meyer.cX, Complexity.Meyer.decode_symBit, Complexity.Meyer.goodF, Complexity.Meyer.sv, Complexity.Meyer.symBit, Complexity.Meyer.tvOf, Complexity.Meyer.xinV, Complexity.Meyer.xin_eq, Turing.MultiTapeTM.runFrom_zero}
\difficulty{5}
Let $M$ be a finite Turing machine with Boolean tape alphabet, let $x$ be a finite
Boolean input word, let $W\in\mathbb{N}$, and let
$xb:\mathbb{N}\to\{\mathrm{false},\mathrm{true}\}$. Suppose every local checker
condition holds for every time $t<2^W$, address $a<2^W$, position $j\le |x|+1$, and
checker. Also suppose $xb(j)=x_j$ for every $j<|x|$. Then at time zero the encoded state
and output register agree with the machine's initial configuration; every encoded
work-tape symbol at offset $d$ agrees with the initial work tape whenever $|d|\le
2^W-1$; and every encoded input symbol at offset $d$ agrees with the initial input tape
whenever $|d|\le |x|+1$.
\end{theorem}

\begin{theorem}[Conditional choice for auxiliary counters]
\lean{Complexity.Meyer.avOf_ite}
\label{Complexity.Meyer.avOf_ite}
\leanok
\uses{Complexity.Meyer.AW, Complexity.Meyer.avOf}
\difficulty{1}
For natural numbers $a$ and $j$, a Boolean $s$, and symbols $A,B\in\mathrm{AW}$, let
$v_{a,j}$ assign the values $a$, $a+1$, $0$, $1$, $j$, and $j+1$ to the six symbols
$ac$, $an$, $a0$, $a1$, $aj$, and $aj1$, respectively. Then $v_{a,j}$ of the symbol
chosen as $A$ when $s$ is true and $B$ otherwise equals $v_{a,j}(A)$ when $s$ is true
and $v_{a,j}(B)$ otherwise.
\end{theorem}

\begin{theorem}[One-Step Work Tape Update]
\lean{Complexity.Meyer.work_edge}
\label{Complexity.Meyer.work_edge}
\leanok
\uses{Complexity.Meyer.AllGood, Complexity.Meyer.Correct, Complexity.Meyer.NumVal, Complexity.Meyer.SS, Complexity.Meyer.aE, Complexity.Meyer.avOf, Complexity.Meyer.avOf_ite, Complexity.Meyer.cX, Complexity.Meyer.claims_of_correct, Complexity.Meyer.edge_E0, Complexity.Meyer.edge_E1, Complexity.Meyer.edge_c1, Complexity.Meyer.edge_c2, Complexity.Meyer.goodF, Complexity.Meyer.reads, Complexity.Meyer.sE, Complexity.Meyer.sv, Complexity.Meyer.sv_sgn, Complexity.Meyer.symvA_NumVal, Complexity.Meyer.tvOf, Complexity.Meyer.workOK, Complexity.Meyer.wsym, Turing.MultiTapeTM.initCfg, Turing.MultiTapeTM.runFrom}
\difficulty{6}
Let $M$ be a finite-tape Turing machine over Booleans, let $x$ be a Boolean input, let
$W$ be a natural number, and let $x_b$ be Boolean-valued data on natural numbers.
Suppose the encoding satisfies every check for times and addresses below $2^W$, indices
through $|x|+1$, and all checks, and is correct at time $t$. Assume $t+1<2^W$, let $h$
be a work-tape index, and let $e$ be an integer with $a_E(e)+1<2^W$. Write $c$ for the
machine configuration after $t$ steps, and let $\rho$ record the symbols read by its
input and work-tape heads. If $c$ is halting, the encoded symbol at offset $e$ is
unchanged at the next time. Otherwise, let $q$ be its state, $m$ the movement of tape
$h$ prescribed by the transition from $q$ on reads $\rho$, and $w$ its prescribed
optional write. The next encoded symbol at offset $e$ equals the symbol selected by
applying $w$ at the scanned cell and moving the reference offset by $m$: for $m=1$, it
is the resulting written-or-retained symbol when $e+1=0$, and otherwise the encoded
value at time $t$ and offset $e+1$; for $m=-1$, the same rule uses offset $e+1$ on the
left side and offset $e$ on the right, with the written-or-retained symbol when $e=0$;
for $m=0$, it is the written-or-retained symbol when $e=0$ and otherwise the encoded
value at time $t$ and offset $e$.
\end{theorem}

\begin{theorem}[Work-Tape Rule Congruence]
\lean{Complexity.Meyer.ruleW_congr}
\label{Complexity.Meyer.ruleW_congr}
\leanok
\uses{Complexity.Meyer.Rd, Complexity.Meyer.ruleW}
\difficulty{4}
Let $M$ be a finite Turing machine with state set $S$ and $k$ work tapes, let
$q\in\mathrm{Option}(S)$, let $\rho$ consist of an optional Boolean input symbol and the
symbols read on the work tapes, choose a work-tape index $h<k$, and let
$d\in\mathbb{Z}$. If $f,g:\mathbb{Z}\to\mathrm{Option}(\mathrm{Bool})$ satisfy
$f(e)=g(e)$ for every integer $e$ whose absolute value is at most $|d|+1$, then the
work-tape update rule at offset $d$ has the same value when computed using $f$ as when
computed using $g$.
\end{theorem}

\begin{theorem}[Work Tape Update Rule]
\lean{Complexity.Meyer.work_rule}
\label{Complexity.Meyer.work_rule}
\leanok
\uses{Complexity.Meyer.AllGood, Complexity.Meyer.Correct, Complexity.Meyer.aE_add_one, Complexity.Meyer.cX, Complexity.Meyer.reads, Complexity.Meyer.ruleW, Complexity.Meyer.work_edge, Turing.MultiTapeTM.initCfg, Turing.MultiTapeTM.runFrom}
\difficulty{5}
Let $M$ be a finite-tape Turing machine over Booleans, let $x$ be a Boolean input, let
$W$ be a natural number, and let $x_b$ be Boolean-valued data on natural numbers.
Suppose the encoding satisfies every check for times and addresses below $2^W$, indices
through $|x|+1$, and all checks, and is correct at time $t$. If $t+1<2^W$, $h$ is a
work-tape index, and $d$ is an integer with $|d|+t+1\le 2^W-1$, then the encoded symbol
on tape $h$ at time $t+1$ and offset $d$ equals the result of applying the machine's
work-tape update rule to the state and reads in its configuration after $t$ steps, using
the time-$t$ encoded tape values as the old tape contents.
\end{theorem}

\begin{theorem}[Correctness of the Work Tape Encoding]
\lean{Complexity.Meyer.work_correct}
\label{Complexity.Meyer.work_correct}
\leanok
\uses{Complexity.Meyer.AllGood, Complexity.Meyer.Correct, Complexity.Meyer.RW, Complexity.Meyer.RW_step, Complexity.Meyer.cX, Complexity.Meyer.ruleW_congr, Complexity.Meyer.work_rule, Turing.MultiTapeTM.initCfg, Turing.MultiTapeTM.runFrom, Turing.MultiTapeTM.runFrom_succ_eq_step'}
\difficulty{3}
Let $M$ be a finite-tape Turing machine over Booleans, let $x$ be a Boolean input, let
$W$ be a natural number, and let $x_b$ be a Boolean-valued function on natural numbers.
Suppose the encoding is good for every time and address below $2^W$, every index $j$ up
to $|x|+1$, and every check, with input-bit data supplied by $x_b$. If $t+1<2^W$, the
encoding is correct at time $t$, $h$ is a work-tape index, and $d$ is an integer
satisfying $|d|+t+1\le 2^W-1$, then the encoded symbol on tape $h$ at time $t+1$ and
offset $d$ equals the symbol actually stored there in the machine's configuration after
$t+1$ steps on input $x$.
\end{theorem}
